\begin{enunciado}{\ejercicio[Diferencia Mínima]}

  \textit{Diferencia Mínima}

     Se tienen dos arreglos de $n$ naturales $A$ y $B$. $A$ está ordenado de manera creciente y B está ordenado de
    manera decreciente. Ningún valor aparece más de una vez en el mismo arreglo. Para cada posición $i$ consideramos
    la diferencia absoluta entre los valores de ambos arreglos $|A[i]-B[i]|$. Se desea buscar el mínimo valor posible
    de dicha cuenta. Por ejemplo, si los arreglos son $A = [1, 2, 3, 4]$ y $B = [6, 4, 2, 1]$ los valores de las diferencias son
    5, 2, 1, 3 y el resultado es 1.

  \begin{enumerate}[label=\arabic*.]
    \item Implementar la función m \textit{minDif}, que tome a $A$ y $B$ y resuelva el problema planteado.

    \item Calcular y justificar la complejidad del algoritmo propuesto. La solución debe ser de tiempo $O(\log n)$,
        dónde $n = tam(A) = tam(B)$.     
  \end{enumerate}
\end{enunciado}

\begin{enumerate}[label=\arabic*)]
  \item Hecho junto al 2.

  \item Tenemos que $A$ es creciente y $B$ decreciente. Notemos lo siguiente:
                \begin{itemize}
                    \item Si $A[i] \geq B[i]$, tenemos que $|A[i]-B[i]| = A[i] - B[i]$, además, 
                            usando que $A[j] > A[i]$ y $B[i] > B[j]$ para $j > i$, obtenemos que
                            $$
                            |A[j]-B[j]| = A[j]-B[j] > A[i] - B[i] = |A[i]-B[i]|
                            $$

                            Esto nos dice que si estamos parados en $i$ y $A[i] \geq B[i]$ , el mínimo debe estar en una posición menor o igual a $i$.
                    \item Si $B[i] > A[i]$, tenemos que $|A[i]-B[i]| = B[i]-A[i]$, además,
                            usando que $A[j] < A[i]$ y $B[i] < B[j]$ para $j < i$, obtenemos que
                            $$
                            |B[j]-A[j]| = B[j]-A[j] > B[i] - A[i] = |A[i]-B[i]|
                            $$
                            Esto nos dice que si estamos parados en $i$ y $A[i] < B[i]$ , el mínimo debe estar en una posición mayor o igual a $i$.
                \end{itemize}
        Conclusión: sale una búsqueda binaria con fritas.

        \lstset{
          mathescape=true,
          emph={[1]funcion, return, ret, retorno},
          emph={[2]for,each, si, sino, if, else, true, false, not},
          emph={[3]minDif},
          emphstyle={[1]\color{violet}\it},
          emphstyle={[2]\color{red}\it},
          emphstyle={[3]\color{OliveGreen}\it},
          morecomment={[l]{//}},
          commentstyle={\color{gray}\it\footnotesize}
        }
        \begin{tcolorbox}
          \begin{lstlisting}
funcion minDif(A, B, izq, der)          
    si izq == der                          
      ret |A[izq]-B[izq]|                           
          
    si der - izq == 1                 // Cuando quedan solo 2 devolvemos el minimo
      ret minimo(|A[izq]-B[izq], A[der]-B[der])

    medio = (izq+der)/2

    si A[medio] >= B[medio]
      ret minDif(A, B, izq, medio)
    sino
      ret minDif(A, B, medio, der)
\end{lstlisting}
        \end{tcolorbox}

Asi, tenemos que $T(n) = O(\log n)$ en el peor caso.
\end{enumerate}